\section{RFD1 versus RFD3: nearly identical pass rates}\label{sec:rfd}
With the metrics reproduced, we compared the two backbone generators head to head. We ran
the same DP3 epitopes through RFdiffusion1 (Lawson's design set) and RFdiffusion3, each
followed by ProteinMPNN sequence design and AlphaFold3 validation, and scored both against
the identical four filters. In aggregate the two are an even trade: RFD1+MPNN passes
\num{840} of \num{151232} designs ($0.56\%$) and RFD3+MPNN passes \num{751} of \num{150948}
($0.50\%$), which is nearly identical. The per-metric distributions are broadly comparable
as well, with RFD3 placing somewhat more mass at low epitope and overall RMSD and a
correspondingly heavier high-PAE tail (Fig.~\ref{fig:rfd}). We read this as license to
standardize on RFD3 going forward: the newer all-atom model reaches the same yield as the
established RFD1 pipeline on this task, so the rest of this log uses RFD3+MPNN unless noted.

\begin{figure}[h]\centering
\figorbox{rfd1_vs_rfd3.png}{1.0}
\caption{RFD1+MPNN (blue, filled) versus RFD3+MPNN (red, outline) on the same DP3 epitopes,
across epitope RMSD, overall RMSD, mean PAE, and AF3 clashing residues, density-normalized.
The four-filter pass rates (in the title) are nearly identical. Regenerated by
\texttt{episcaf\_analysis/viz/plot\_rfd1\_vs\_rfd3.py}.}
\label{fig:rfd}
\end{figure}
